\section{Worked examples: wheat in India}\label{sec:worked}
The three examples follow one pair, wheat in India, from saved inputs to the reported numbers. In each case the
code (\file{17\_compare\_runs.py}, \file{26\_report\_tables.py}) recomputes the result by hand and asserts that it
equals the pipeline value.

\subsection{From raw components to F4}
Table~\ref{tab:worked} shows the present-day calculation. It also explains why scores roughly halved between the
four-crop and nine-crop runs.
\begin{itemize}
\item \textbf{Same physical inputs.} The heat level of wheat in India is \exNineLheatc\,\si{\celsius} of $\Teff$
above \SI{25}{\celsius} in both runs.
\item \textbf{Different maxima.} In the nine-crop run the heat maximum is set by common bean in India
(\exNinemaxLheat\,\si{\celsius}), so wheat's heat exposure falls from \exFourEheat\ to \exNineEheat.
\item \textbf{Result.} The water term dominates:
$\Fd = \tfrac12(\exNineEheat\times\exNineSheat + \exNineEwater\times\exNineSwater)\times\exNineInorm = \exNineFFourpipeline$.
\end{itemize}

% worked example: wheat in India, from raw components to F4 (values from worked_example.csv via macros)
\begin{table}[t]
\centering\small
\caption{Worked example: wheat in India. The physical inputs are identical in both runs; only the
normalisation maxima differ. Values from \texttt{worked\_example.csv}; the hand calculation
$\tfrac12(E_TS_T+E_PS_P)\tilde I$ equals the pipeline value (checked by an assertion in the code).}\label{tab:worked}
\begin{tabular}{lrr}
\toprule
 & four-crop run & nine-crop run\\
\midrule
Heat level $L_{\text{heat}}$ (\si{\celsius} of $\Teff$ above \SI{25}{\celsius}) & \exFourLheatc & \exNineLheatc\\
\quad maximum over pairs & \exFourmaxLheat & \exNinemaxLheat\\
\quad $E_{\text{heat}}$ & \exFourEheat & \exNineEheat\\
Water level $L_{\text{water}}$ (rainfed share $\times$ deficit) & \exFourLwater & \exNineLwater\\
\quad maximum over pairs & \exFourmaxLwater & \exNinemaxLwater\\
\quad $E_{\text{water}}$ & \exFourEwater & \exNineEwater\\
Heat slope $s_T$ (/\si{\celsius}) and maximum over crops & \exFoursTpc\ / \exFourmaxsT & \exNinesTpc\ / \exNinemaxsT\\
\quad $S_T$ & \exFourSheat & \exNineSheat\\
$K_y$ and maximum over crops & \exFourKy\ / \exFourmaxKy & \exNineKy\ / \exNinemaxKy\\
\quad $S_P$ & \exFourSwater & \exNineSwater\\
Importance $I$ and maximum & \exFourI\ / \exFourmaxI & \exNineI\ / \exNinemaxI\\
\quad $\tilde I$ & \exFourInorm & \exNineInorm\\
\midrule
\Fd & \exFourFFourpipeline & \exNineFFourpipeline\\
\bottomrule
\end{tabular}
\end{table}


\subsection{A standard 10\% loss of the wheat crop}
In 2021--23, India produced \wsWheatIndProd\,kt of wheat against a domestic supply of \wsWheatIndDom\,kt and
exports of \wsWheatIndExp\,kt, so SSR = \wsSSR\ and $\text{SSR}^{*}=1$.

\paragraph{Upper bound (no market response).} A 10\% yield loss removes $f^{\text{up}} = 0.10\times1 =
\wsFrac$ of wheat's contribution to supply.
\begin{itemize}
\item \textbf{Iron.} Wheat supplies \wsCIron\ of India's \wsBIron\,mg iron per person per day, so iron supply
changes by $-0.10\times\wsCIron/\wsBIron$, that is \wsDIron\%. Iron adequacy moves from
$\min(1, \wsBIron/\wsRIron) = \wsNARzeroIron$ to \wsNARoneIron.
\item \textbf{Energy and zinc.} Energy changes by \wsDKcal\% and zinc by \wsDZinc\%. Both stay above their
requirements, so their NAR stays 1.
\item \textbf{MAR.} The MAR falls from \wsMARzero\ to \wsMARone, a change of \wsDMAR\ points. The whole change
comes from iron.
\end{itemize}

\paragraph{Lower bound (exports absorb first).} Exports absorb part of the loss, so $f^{\text{low}} = \wsLower$ and
the MAR change is \wsDMARlower\ points.

\paragraph{Reading.} The MAR cap hides the energy and zinc losses: they are real but leave supply above the
requirement. This is why the nutrient changes are reported beside MAR.

\subsection{The conditional future yield change}
\paragraph{Exposure.} From 1971--2000 to 2041--2060 under MRI-ESM2-0 SSP2-4.5, the flowering-window $\Teff$ of
Indian wheat rises by \wfdTeff\,\si{\celsius}.
\begin{itemize}
\item The heat level rises from \wfHheatHist\ to \wfHheatFut\,\si{\celsius}.
\item The water level barely moves (\wfHwaterHist\ to \wfHwaterFut).
\end{itemize}

\paragraph{Index.} With the fixed reference, $E_{\text{heat}}$ rises from \wfEheatHist\ to \wfEheatFut\ and
\Fd\ from \wfFFourHist\ to \wfFFourFut. Of that change, \wfdFfourHeat\ is the heat part and \wfdFfourWater\ the water
part (Eq.~\ref{eq:decomp}).

\paragraph{Yield.} In physical units (Eq.~\ref{eq:loss}), with $s_T=\wfSlope$ and $K_y=\wfKy$:
\begin{itemize}
\item the period-mean heat loss rises from \wfHeatLossHist\ to \wfHeatLossFut;
\item the water loss rises from \wfWaterLossHist\ to \wfWaterLossFut;
\item the combined loss rises from \wfLossHist\ to \wfLossFut;
\item hence $\Delta Y = -100\,(\wfLossFut-\wfLossHist)/(1-\wfLossHist) = \wfdY\%$.
\end{itemize}
The empirical coefficients give \wfdYemp\%.

\paragraph{Nutrition.} Propagated through Eq.~\ref{eq:shock} (upper bound), the change reduces India's iron supply
by \wfDIron\% and its energy supply by \wfDKcal\%, and lowers MAR by \wfDMAR\ points.

\paragraph{Reading.} The index change (\wfFFourHist\ to \wfFFourFut) and the yield change (\wfdY\%) are different
quantities. Only the second has physical units, and it is conditional on the literature damage functions.
